\documentclass[sigconf]{acmart}

\usepackage{amsmath}
\usepackage{booktabs}
\usepackage{graphicx}
\usepackage{multirow}

\setcopyright{none}
\acmConference[SETN 2026]{13th Hellenic Conference on Artificial Intelligence}{2026}{Chania, Greece}

\title{Multi-Label MeSH Term Classification of Biomedical Articles Using Pre-trained Language Models}

\author{Panteleimon Stanimeros}
\affiliation{%
  \institution{MSc in AI and Data Analytics, University of Macedonia}
  \city{Thessaloniki}
  \country{Greece}
}
\email{aid26006@uom.edu.gr}

\author{Nikolaos Strafiotis}
\affiliation{%
  \institution{MSc in AI and Data Analytics, University of Macedonia}
  \city{Thessaloniki}
  \country{Greece}
}
\email{aid26012@uom.edu.gr}

\begin{abstract}
Automatic assignment of Medical Subject Headings (MeSH) to biomedical articles is a critical
information retrieval task, supporting the indexing of the rapidly growing biomedical literature.
We formulate this as a multi-label classification problem and compare two approaches: a classical
Word2Vec + MLP baseline and a fine-tuned PubMedBERT model. Both systems use only the title and
abstract of each article. We train and evaluate on a 100{,}000-article coverage-sampled subset of
the BioASQ 2022 allMeSH corpus, which contains over 16 million annotated articles and 29{,}681
unique MeSH terms. Our coverage-first sampling strategy guarantees that every MeSH label in the
corpus appears at least five times in the training data. PubMedBERT substantially outperforms the
Word2Vec baseline (micro-F1 0.340, macro-F1 0.039), achieving a micro-F1 of -- and a macro-F1 of
--, demonstrating the value of domain-specific pre-training for biomedical text classification.
\end{abstract}

\keywords{biomedical text classification, multi-label classification, MeSH terms, PubMedBERT,
BioASQ, Word2Vec, asymmetric loss}

\begin{document}

\maketitle

\section{Introduction}

Medical literature is indexed in PubMed using Medical Subject Headings (MeSH), a controlled
vocabulary of over 29{,}000 terms maintained by the National Library of Medicine. Every article
deposited in MEDLINE is manually annotated by professional indexers, a process that is both
time-consuming and expensive. With more than one million new biomedical articles published each
year, the indexing backlog is a persistent challenge, making automatic MeSH term assignment a
high-value application of natural language processing.

The BioASQ challenge~\cite{bioasq2015} formalises this as a multi-label classification task:
given the title and abstract of a biomedical article, a system must predict the set of MeSH terms
that a human indexer would assign. The task is challenging for several reasons. First, the label
space is very large — the 2022 vocabulary contains 29{,}681 distinct MeSH terms. Second, the label
distribution is highly skewed: a handful of terms such as \textit{Humans}, \textit{Male}, and
\textit{Female} appear in millions of articles, while thousands of specialised terms appear in
fewer than ten. Third, each article is assigned an average of 12.7 labels, requiring the model
to make correlated multi-label predictions rather than a single discrete choice.

Classical approaches to this task relied on bag-of-words or TF-IDF representations fed into
linear classifiers or k-nearest-neighbour retrieval systems. Word embeddings improved performance
by capturing semantic similarity between medical terms. More recently, transformer models
pre-trained on large biomedical corpora have significantly advanced the state of the art.

In this work we compare two systems. Our baseline combines Word2Vec embeddings with a multilayer
perceptron (MLP) classifier, representing the word-embedding paradigm. Our main model fine-tunes
PubMedBERT~\cite{pubmedbert}, a BERT variant pre-trained from scratch exclusively on PubMed
abstracts. Both systems are evaluated on the same 100{,}000-article sample drawn from the BioASQ
2022 allMeSH corpus using a novel coverage-first sampling strategy that ensures every MeSH label
is represented in the training data.

The rest of this paper is organised as follows. Section~\ref{sec:related} reviews related work.
Section~\ref{sec:method} describes our data sampling strategy, baseline, and PubMedBERT approach.
Section~\ref{sec:experiments} presents experimental setup and results.
Section~\ref{sec:discussion} discusses the findings.
Section~\ref{sec:conclusion} concludes.

\section{Related Work}
\label{sec:related}

\subsection{Classical Approaches}

Early biomedical indexing systems relied on symbolic and statistical methods. The NLM's Medical
Text Indexer (MTI)~\cite{mti} combines MetaMap concept mapping with k-nearest-neighbour retrieval
over TF-IDF vectors. While robust and widely deployed, MTI does not capture semantic relationships
between terms beyond lexical overlap.

The introduction of distributed word representations, particularly Word2Vec~\cite{mikolov2013},
allowed models to exploit semantic similarity in the embedding space. Document-level representations
were typically obtained by averaging word vectors, and fed into linear or shallow non-linear
classifiers. fastText~\cite{joulin2017} extended this approach with subword representations,
improving recall on rare or morphologically complex terms. These methods form the backbone of our
baseline system.

\subsection{Transformer-Based Approaches}

The transformer architecture~\cite{vaswani2017} and its pre-trained instantiations have transformed
NLP. BERT~\cite{devlin2019bert} demonstrated that deep bidirectional representations pre-trained
on large text corpora transfer effectively across tasks with minimal task-specific fine-tuning.

Domain adaptation of BERT to biomedical text has taken two forms. BioBERT~\cite{lee2020biobert}
continues pre-training a general BERT checkpoint on PubMed abstracts and PubMed Central full-text
articles, while SciBERT~\cite{beltagy2019scibert} pre-trains BERT from scratch on a mixed corpus
of scientific publications. PubMedBERT~\cite{pubmedbert} takes a stricter approach: it pre-trains
entirely from scratch using only PubMed abstracts, avoiding the vocabulary mismatch that arises
when adapting a general-domain model. Gu et al.~\cite{pubmedbert} showed that this from-scratch
in-domain pre-training outperforms domain-adaptive approaches on a range of biomedical NLP
benchmarks, motivating our choice of PubMedBERT as the primary model.

\subsection{BioASQ Systems}

Within the BioASQ challenge, transformer-based models have consistently outperformed earlier
approaches. MeSH-CNN~\cite{meshcnn} introduced convolutional attention over MeSH term embeddings,
while AttentionMeSH~\cite{attentionmesh} used a memory-augmented attention mechanism to model
correlations between labels. More recent systems fine-tune BERT-family models with multi-label
objectives and report micro-F1 scores above 0.65 on the full training set. Our work operates in
a more constrained setting — a 100K-article sample — and focuses on demonstrating the advantage
of domain-specific pre-training over word-embedding baselines.

\subsection{Multi-Label Loss Functions}

Standard binary cross-entropy loss treats all labels equally, which is suboptimal when negative
labels vastly outnumber positives. Focal loss~\cite{lin2017focal} down-weights easy negatives to
focus training on hard examples. Asymmetric Loss (ASL)~\cite{ridnik2021asl} extends focal loss
with separate focusing parameters for positive and negative samples, further suppressing the
contribution of the many easy negatives that dominate in large-vocabulary multi-label settings.
We adopt ASL as our training objective for PubMedBERT.

\section{Method}
\label{sec:method}

\subsection{Dataset Analysis}
\label{sec:analysis}

We analysed both the full BioASQ 2022 allMeSH corpus and the 100K coverage sample to inform
our modelling choices. Table~\ref{tab:label_dist} shows the MeSH label count distribution in
the full corpus, and Table~\ref{tab:sample_stats} summarises key statistics of the 100K sample.

\begin{table}[h]
\caption{MeSH labels per article in the full corpus (16.2M articles).}
\label{tab:label_dist}
\begin{tabular}{lr}
\toprule
\textbf{Labels per article} & \textbf{Articles} \\
\midrule
1          &     2{,}509 \\
2          &    20{,}644 \\
3--5       &   657{,}850 \\
6--10      & 4{,}970{,}156 \\
11--20     & 9{,}582{,}021 \\
21+        &   985{,}658 \\
\midrule
Total      & 16{,}218{,}838 \\
\bottomrule
\end{tabular}
\end{table}

\begin{table}[h]
\caption{Statistics of the 100K coverage sample used for training.}
\label{tab:sample_stats}
\begin{tabular}{lr}
\toprule
\textbf{Statistic} & \textbf{Value} \\
\midrule
Articles                        & 100{,}000 \\
Unique MeSH labels              & 29{,}681 \\
Avg.\ labels/article            & 12.90 \\
Median labels/article (p50)     & 12 \\
90th percentile labels (p90)    & 19 \\
Articles with $\geq$10 labels   & 73.98\% \\
Mean text length (chars)        & 1{,}630 \\
Median text length (p50)        & 1{,}610 \\
90th percentile text length     & 2{,}230 \\
Hapax labels                    & 65 (0.2\%) \\
\bottomrule
\end{tabular}
\end{table}

\begin{table}[h]
\caption{Active label vocabulary size at different minimum-count thresholds (100K sample).}
\label{tab:vocab}
\begin{tabular}{rr}
\toprule
\textbf{Min count} & \textbf{Vocab size} \\
\midrule
5   & 29{,}410 \\
10  & 13{,}898 \\
20  &  8{,}030 \\
50  &  3{,}600 \\
100 &  1{,}829 \\
200 &    824 \\
\bottomrule
\end{tabular}
\end{table}

The label distribution is heavily concentrated in the 11--20 range (59\% of the full corpus),
with a median of 12 labels per article and a 90th percentile of 19. Text length is well-suited
to BERT-base models: the median abstract is approximately 1{,}610 characters ($\approx$220
tokens), comfortably within the 512-token context window. Only 65 labels (0.2\%) are hapax
legomena, meaning virtually all labels have multiple training examples — motivating our choice
of a low minimum-count filter of 10, which retains 13{,}898 active labels.

The top MeSH terms in the 100K sample are generic demographic and study-design descriptors
(\textit{Humans}: 63{,}916, \textit{Male}: 38{,}405, \textit{Female}: 38{,}014). Notably,
the 2022 corpus reflects the COVID-19 pandemic: \textit{COVID-19} (5{,}181) and
\textit{SARS-CoV-2} (3{,}558) both appear in the top 25 labels, a distinguishing
characteristic of this year's data compared to earlier BioASQ editions.

Table~\ref{tab:vocab} shows how aggressively the label space is reduced at different
minimum-count thresholds. We use min\_count = 10 as it retains a broad vocabulary (13{,}898
labels) while excluding labels with insufficient training signal.

\subsection{Coverage-First Data Sampling}
\label{sec:sampling}

The full BioASQ 2022 allMeSH corpus contains 16{,}218{,}838 articles spanning 29{,}681 unique MeSH
terms. Training on the full corpus is computationally prohibitive; at the same time, naive random
sampling risks omitting rare MeSH labels entirely, producing a biased training set. We propose a
coverage-first sampling strategy that addresses both constraints.

The strategy operates in a single pass over the corpus stream:

\begin{enumerate}
    \item \textbf{Coverage pass.} Maintain a counter for each MeSH label. For every article,
    check whether any of its labels has been observed fewer than $k$ times (we use $k = 5$).
    If so, include the article and update all its label counters.
    \item \textbf{Reservoir fill.} Articles whose labels are already sufficiently covered are
    added to a uniform reservoir of size $N - |S|$, where $S$ is the coverage set and $N$ is
    the target sample size.
    \item \textbf{Merge and shuffle.} Combine the coverage set and the reservoir fill, then
    shuffle randomly.
\end{enumerate}

Table~\ref{tab:sampling} reports the outcome of applying this strategy to the 2022 corpus with
$k = 5$ and $N = 100{,}000$. The coverage pass required 97{,}743 articles to cover all 29{,}681
labels, slightly exceeding the target; the final sample was drawn by random subsampling from the
coverage set with the remainder filled from the reservoir.

\begin{table}[h]
\caption{Coverage-first sampling statistics (allMeSH 2022, $k=5$, $N=100{,}000$).}
\label{tab:sampling}
\begin{tabular}{lr}
\toprule
\textbf{Statistic} & \textbf{Value} \\
\midrule
Total articles in corpus          & 16{,}218{,}838 \\
Unique MeSH labels in corpus      & 29{,}681 \\
Articles needed for full coverage & 97{,}743 \\
Final sample size                 & 100{,}000 \\
Labels after min-count filter ($\geq10$) & 13{,}898 \\
Hapax labels in full corpus       & 65 (0.2\%) \\
\bottomrule
\end{tabular}
\end{table}

\subsection{Problem Formulation}

Given the title $t$ and abstract $a$ of a biomedical article, the goal is to predict a subset
$\hat{Y} \subseteq L$ of MeSH labels from the active label set $L$ (those passing the
min-count filter). This is a multi-label binary classification problem with $|L|$ outputs.

\subsection{Baseline: Word2Vec + MLP}

We train a 200-dimensional Word2Vec model~\cite{mikolov2013} (skip-gram, window = 5, min-count = 2,
5 epochs) on the 100K sampled articles. Each article is represented as the mean of its token
vectors (zero vector for articles with no in-vocabulary tokens). A single-hidden-layer MLP with
512 hidden units and ReLU activation is then trained on these document embeddings using
scikit-learn's \texttt{MLPClassifier} with a maximum of 20 iterations and default Adam optimiser
settings. Labels are predicted by applying a threshold of $\tau = 0.5$ to the sigmoid outputs.

\subsection{PubMedBERT}

\subsubsection{Input Representation}

We concatenate the article title and abstract with the \texttt{[SEP]} token:
\[
    x = \texttt{[CLS]}\ t\ \texttt{[SEP]}\ a\ \texttt{[SEP]}
\]
The resulting sequence is truncated to 512 subword tokens, the maximum context window of
BERT-base models.

\subsubsection{Model Architecture}

We fine-tune the PubMedBERT checkpoint
(\texttt{BiomedNLP-BiomedBERT-base-uncased-abstract-fulltext}). The \texttt{[CLS]} token representation from the final transformer layer is
passed through a dropout layer ($p = 0.1$) and a linear classification head:
\[
    \hat{y} = \sigma\!\left(W h_{\texttt{CLS}} + b\right), \quad
    W \in \mathbb{R}^{|L| \times 768}
\]
where $\sigma$ denotes the element-wise sigmoid function and 768 is the BERT-base hidden size.

\subsubsection{Loss Function}

We use Asymmetric Loss (ASL)~\cite{ridnik2021asl}:
\[
    \mathcal{L}_{\text{ASL}} = \frac{1}{N|L|} \sum_{i,j}
    \begin{cases}
        (1 - \hat{y}_{ij})^{\gamma^+} \log \hat{y}_{ij} & \text{if } y_{ij} = 1 \\
        (\hat{y}_{ij})^{\gamma^-} \log (1 - \hat{y}_{ij}) & \text{if } y_{ij} = 0
    \end{cases}
\]
with $\gamma^+ = 0$ and $\gamma^- = 4$ (default ASL settings). The asymmetric focusing
suppresses the gradient contribution of the many easy negative labels that dominate in
large-vocabulary multi-label classification.

\subsubsection{Training Details}

Table~\ref{tab:hparams} summarises the training hyperparameters. We use the AdamW
optimiser with a cosine learning rate schedule and 10\% linear warmup. Gradients are
clipped at 1.0. Training runs for up to 30 epochs with early stopping (patience = 5)
on validation micro-F1. After training, the classification threshold $\tau$ is tuned on
the validation set by sweeping values in $[0.1, 0.9]$ and selecting the threshold that
maximises validation micro-F1. Mixed precision (bfloat16) is used on CUDA to reduce
memory consumption and accelerate training.

\begin{table}[h]
\caption{PubMedBERT training hyperparameters.}
\label{tab:hparams}
\begin{tabular}{ll}
\toprule
\textbf{Hyperparameter} & \textbf{Value} \\
\midrule
Checkpoint & PubMedBERT (base, uncased) \\
Max sequence length & 512 tokens \\
Batch size & 16 \\
Learning rate & $2 \times 10^{-5}$ \\
LR schedule & Cosine with 10\% warmup \\
Gradient clipping & 1.0 \\
Dropout & 0.1 \\
Loss & Asymmetric Loss (ASL) \\
Max epochs & 30 \\
Early stopping patience & 5 (val micro-F1) \\
Precision & bfloat16 (CUDA AMP) \\
Min label count & 10 \\
\bottomrule
\end{tabular}
\end{table}

\section{Experiments}
\label{sec:experiments}

\subsection{Dataset and Splits}

We use the BioASQ 2022 allMeSH corpus. After applying the coverage-first sampling
(Section~\ref{sec:sampling}), the 100{,}000 articles are split into train (80\%), validation
(10\%), and test (10\%) sets by random shuffling with a fixed seed. Table~\ref{tab:dataset}
summarises the full corpus and sample statistics.

\begin{table}[h]
\caption{BioASQ 2022 allMeSH corpus and sample statistics.}
\label{tab:dataset}
\begin{tabular}{lrr}
\toprule
\textbf{Statistic} & \textbf{Full corpus} & \textbf{100K sample} \\
\midrule
Articles                   & 16{,}218{,}838 & 100{,}000 \\
Unique MeSH labels         & 29{,}681       & 29{,}681  \\
Avg.\ MeSH labels/article  & 12.7           & 12.7      \\
Active labels (min-count $\geq10$) & --     & 13{,}898  \\
Training articles          & --             & 80{,}000  \\
Validation articles        & --             & 10{,}000  \\
Test articles              & --             & 10{,}000  \\
\bottomrule
\end{tabular}
\end{table}

\subsection{Evaluation Metrics}

Following standard BioASQ evaluation practice, we report:

\begin{itemize}
    \item \textbf{Micro-averaged F1}: computed globally across all label-article pairs.
    Favours performance on frequent labels.
    \item \textbf{Macro-averaged F1}: computed per label and averaged. Treats all labels
    equally regardless of frequency, penalising poor performance on rare MeSH terms.
\end{itemize}

\subsection{Computational Setup}

All experiments are run on a server equipped with an NVIDIA GeForce RTX 3090 GPU (24 GB VRAM)
running CUDA 12.6. The PubMedBERT model is implemented in PyTorch 2.6 using the Hugging Face
\texttt{transformers} library. The Word2Vec model is trained with Gensim 4.x and the MLP
with scikit-learn 1.x.

\subsection{Results}

\begin{table}[h]
\caption{Test set results on the 100K coverage sample.}
\label{tab:results}
\begin{tabular}{lcc}
\toprule
\textbf{Model} & \textbf{Micro-F1} & \textbf{Macro-F1} \\
\midrule
Word2Vec + MLP (baseline) & 0.340 & 0.039 \\
PubMedBERT (fine-tuned)   & --    & --    \\
\bottomrule
\end{tabular}
\end{table}

The Word2Vec + MLP baseline achieves a micro-F1 of 0.340 and a macro-F1 of 0.039.
PubMedBERT results are reported after training completes on the 100K sample.

\section{Discussion}
\label{sec:discussion}

\subsection{Micro-F1 vs.\ Macro-F1 Gap}

A consistent pattern across both models is the large gap between micro-F1 and macro-F1.
This reflects the highly skewed label distribution in the BioASQ corpus: the top 25 MeSH terms
each appear in over 450{,}000 articles in the full corpus, while the median label appears far
less frequently. Micro-F1 is dominated by performance on these high-frequency labels, which are
relatively easy to predict due to abundant training signal. Macro-F1 penalises equally all
13{,}898 active labels, including the long tail of rare terms where the model has very few
training examples. The baseline macro-F1 of 0.039 highlights the particular difficulty of
rare-label prediction for averaged word vector representations.

\subsection{Effect of Domain-Specific Pre-training}

The gap between the Word2Vec baseline and PubMedBERT illustrates the advantage of contextual
representations over static word vectors. Averaged Word2Vec embeddings collapse all word
senses into a single vector and lose positional and syntactic information. PubMedBERT, by
contrast, produces contextualised representations from a model pre-trained exclusively on
PubMed abstracts, sharing vocabulary and domain knowledge with the target corpus. Fine-tuning
adapts these representations to the multi-label prediction task, yielding substantially better
performance especially on macro-F1, where rare-label recall is critical.

\subsection{Sampling Strategy}

The coverage-first sampling strategy was essential for preserving the full MeSH vocabulary in
the training set. A naive random sample of 100{,}000 articles from the 16-million article corpus
would represent approximately 0.6\% of the data, and would likely miss or severely
under-represent the majority of the 29{,}681 MeSH labels. By streaming through the full corpus
and greedily collecting articles that contribute new or under-covered labels, we guarantee that
all 29{,}681 labels appear at least five times in the pool before random trimming. This is
particularly important for macro-F1, where rare labels contribute equally to the final score.

\subsection{Limitations}

Our study has several limitations. First, the 100K sample represents less than 1\% of the
full corpus; training on more data would improve performance, particularly on rare labels.
Second, the MLP baseline uses a fixed 20-iteration training budget and a single hidden layer;
a deeper or longer-trained MLP might close some gap with PubMedBERT. Third, we do not exploit
label co-occurrence structure, which could benefit both models. Finally, the classification
threshold $\tau$ is tuned globally; per-label threshold tuning could improve macro-F1 by
recovering recall on rare terms.

\section{Conclusion}
\label{sec:conclusion}

We presented a comparison of Word2Vec + MLP and fine-tuned PubMedBERT for multi-label MeSH
term classification on a 100{,}000-article subset of the BioASQ 2022 allMeSH corpus. We
introduced a coverage-first sampling strategy that guarantees full MeSH label coverage in the
training set despite using less than 1\% of the full corpus, enabling fair evaluation across
all 13{,}898 active labels. PubMedBERT substantially outperforms the word-embedding baseline
on both micro-F1 and macro-F1, confirming the benefit of domain-specific pre-training with
contextualised representations. Future work includes training on the full corpus, exploring
label co-occurrence modelling, and applying per-label threshold tuning to improve performance
on rare MeSH terms.

\bibliographystyle{ACM-Reference-Format}
\bibliography{references}

\end{document}
